\thesischapter{III}{METODE PENELITIAN}

\section{Metode Pendekatan}\label{sec:metode-pendekatan}

Penelitian ini menggunakan pendekatan kuantitatif. Penelitian kuantitatif menguji teori secara objektif dengan memeriksa hubungan antarvariabel, dan variabel tersebut diukur dengan instrumen sehingga data berupa angka dapat dianalisis secara statistik \parencite{creswell2018}. Datanya berupa jumlah pelanggaran kaidah dan pelanggaran per birama pada harmonisasi yang dihasilkan model, lalu dibandingkan antarmodel dan antarasal melodi melalui uji statistik. Teori harmoni fungsional dan kaidah Strube digunakan secara deduktif, yaitu kaidah diturunkan menjadi variabel, hipotesis, dan definisi operasional sebelum data dikumpulkan.

Pertanyaan pertama menggambarkan tingkat kepatuhan setiap model pada setiap asal melodi. Pertanyaan kedua membandingkan dua model pada melodi yang sama. Kedua model dipilih peneliti dan tidak diberi perlakuan, sehingga perbandingan ini bersifat komparatif. Pertanyaan ketiga membandingkan melodi dari dua asal yang diberikan peneliti kepada model yang sama. Pengaruh asal melodi diperiksa dengan logika eksperimen, yaitu mencari pengaruh perlakuan tertentu dalam kondisi yang dikendalikan \parencite{sugiyono2019}. Asal melodi tidak dapat diacak pada melodi yang sama, karena setiap melodi memang berasal dari satu sumber. Jadi, yang diperiksa adalah pengaruh asal melodi secara keseluruhan, bukan pengaruh satu ciri melodi tertentu.

Desain penelitian mengikuti analisis Bach Doodle \parencite[798]{huang2019} dengan empat perluasan yang dirangkum dalam Tabel~\ref{tab:posisi-penelitian}. Seluruh pengukuran dilakukan secara otomatis oleh program komputer. Penelitian ini tidak menggunakan penilai, pendengar, maupun partisipan manusia. Contoh partitur pada pembahasan hanya ilustrasi dan tidak diolah sebagai data, sehingga penelitian ini tidak menggunakan metode campuran (\textit{mixed methods}).

Penelitian berlangsung dalam dua tahap. Tahap pertama merumuskan kaidah serta membangun dan memvalidasi instrumen evaluasi otomatis berbasis music21. Tahap kedua menerapkan instrumen tersebut pada harmonisasi kedua model untuk menjawab pertanyaan penelitian.

Perbandingan antarmodel adalah studi komparatif, bukan eksperimen yang terkontrol penuh. DeepBach dan Coconet sama-sama dilatih pada \textit{chorale} Bach, tetapi versi data, representasi, rancangan jaringan, dan prosedur generasinya berbeda. Perbedaan pelanggaran per birama antarmodel dapat mencerminkan gabungan faktor-faktor tersebut \parencite[11--13]{briot2020book}. Pengaruh masing-masing faktor tidak dapat dipisahkan, karena tidak ada versi kedua model yang dilatih dan dijalankan dengan cara yang identik. Temuan berlaku pada model, \textit{checkpoint}, dan melodi yang diuji. \textit{Checkpoint} adalah berkas bobot hasil pelatihan model yang dipakai untuk menghasilkan keluaran.

Penelitian ini bersifat interdisipliner karena memadukan teori musik dengan komputasi, yaitu pengolahan data oleh program komputer. Evaluasi model generatif musik sendiri dapat dilihat dari sudut pandang musikologi, teknik, dan interaksi manusia--komputer \parencite{lerch2025}. Roberts dan Horn menggambarkan gabungan bermusik dan menulis kode sebagai praktik hibrida, tempat kebiasaan dari kedua bidang saling membentuk ulang \parencite{roberts2025}.

Perumusan kaidah dan pemilihan melodi tetap melibatkan keputusan peneliti. Posisi peneliti dibentuk oleh latar belakang dan pengalamannya serta dapat memengaruhi proses penelitian \parencite{holmes2020}. Peneliti adalah mahasiswa musik yang mempelajari harmoni Barat, sehingga berposisi sebagai orang dalam (\textit{insider}) terhadap tradisi teori yang dijadikan acuan. Terhadap kedua model, peneliti berposisi sebagai orang luar (\textit{outsider}) karena tidak terlibat dalam pengembangan dan pelatihannya. Kedua posisi tersebut dapat ditempati bersamaan \parencite{dwyer2009}. Refleksivitas, yaitu pemeriksaan peneliti atas anggapan dan penilaiannya sendiri, juga berguna pada penelitian kuantitatif \parencite{jamieson2023}. Pada penelitian ini, posisi peneliti dapat memengaruhi perumusan kaidah dan pengecualiannya, pemilihan melodi, dan penafsiran hasil. Karena itu, definisi kaidah, kriteria melodi, pengaturan generasi, dan rencana analisis ditetapkan dalam protokol tertulis sebelum harmonisasi utama dihasilkan. Melodi \textit{chorale} Bach dipilih secara acak dengan \textit{seed} (bilangan acak awal) yang dicatat, seluruh latihan Strube yang memenuhi kriteria diikutsertakan, instrumen yang sama diterapkan secara otomatis pada seluruh keluaran, dan contoh partitur untuk pembahasan dipilih secara acak dari kejadian yang ditandai.

Kaidah Strube mewakili norma penulisan empat suara dalam tradisi harmoni tonal Barat. Tradisi Barat juga mendominasi data penelitian pembuatan musik. Mehta et al.\ mencatat bahwa hanya 5,7 persen dari total jam \textit{dataset} musik yang mereka kaji berasal dari genre non-Barat \parencite{mehta2025}. Karena itu, kejadian yang ditandai menunjukkan ketidaksesuaian dengan norma tradisi tersebut, bukan penilaian atas nilai musikal keluaran atau atas tradisi musik lain.

\section{Objek dan Subjek Penelitian}\label{sec:objek-dan-subjek}

Objek penelitian adalah harmonisasi simbolik empat suara yang dihasilkan DeepBach \parencite{hadjeres2017} dan Coconet \parencite{huang2017} untuk melodi sopran yang ditetapkan. Harmonisasi asli Bach untuk melodi \textit{chorale} yang digunakan juga diukur sebagai acuan. DeepBach dijalankan melalui implementasi PyTorch resmi pengembangnya pada versi kode yang dikunci dalam protokol, beserta bobot hasil pelatihan yang disediakan untuk implementasi tersebut. Coconet dijalankan melalui pustaka Magenta.js versi 1.23.1 dengan \textit{checkpoint} \texttt{coconet/bach} yang disediakan Google. Identitas setiap versi kode, paket, dan berkas bobot dicatat bersama data.

Kedua model dipilih karena kode dan bobotnya tersedia secara terbuka, dapat menerima melodi sopran yang ditetapkan lalu menghasilkan alto, tenor, serta bas, dan dilatih pada \textit{chorale} Bach sehingga melodi \textit{chorale} berasal dari jenis musik yang dipelajarinya. NotaGen \parencite{wang2025}, misalnya, tidak dipakai karena hanya menerima petunjuk periode, komponis, dan instrumentasi, bukan melodi yang harus dipertahankan.

Penelitian ini tidak menjadikan manusia sebagai subjek penelitian dan tidak melibatkan manusia sebagai penilai.

\subsection{Populasi}\label{sec:populasi}

Populasi penelitian ini adalah melodi sopran dari dua sumber. Populasi pertama adalah melodi sopran \textit{chorale} Bach dalam korpus music21 \parencite{cuthbert2010}. Daftar \textit{chorale} dalam korpus tersebut memuat 371 nomor menurut katalog Riemenschneider, tetapi beberapa nomor menunjuk berkas yang sama, sehingga daftar itu berisi 350 berkas yang berbeda. Populasi kedua adalah latihan harmonisasi melodi dalam buku Strube edisi 1928 \parencite{strube1928}. Karena proses generasinya mengandung unsur acak, model dapat menghasilkan harmonisasi yang berbeda untuk melodi yang sama setiap kali dijalankan. Harmonisasi yang mungkin dihasilkan untuk satu melodi tidak dapat dicacah secara lengkap, sehingga diwakili oleh beberapa harmonisasi yang dihasilkan dengan \textit{checkpoint} dan pengaturan generasi yang sama.

\subsection{Sampel}\label{sec:sampel}

Unit analisis penelitian ini adalah melodi. Sampel terdiri atas dua kelompok melodi.

\begin{enumerate}
    \item Kelompok Bach, yaitu melodi sopran \textit{chorale} Bach yang diambil secara acak dari populasi pertama.
    \item Kelompok Strube, yaitu seluruh latihan harmonisasi melodi bernomor yang layak dalam buku Strube, pada bab-bab akor dari harmonisasi melodi pertama sampai bab mode minor (hlm. 11--80).
\end{enumerate}

Kedua kelompok memakai kriteria kelayakan yang sama. Melodi harus berupa satu suara dengan birama dan tangga nada yang tertulis, seluruh nilai nadanya dapat dinyatakan dalam satuan not seperenam belas, dan seluruh nadanya berada dalam wilayah masukan sopran kedua model. Kesesuaian dengan kosakata masukan DeepBach dan wilayah nada Coconet diperiksa oleh program sebelum generasi, tanpa menjalankan model. \textit{Chorale} Bach juga harus memiliki tepat empat suara ketika dibaca dalam satuan not seperenam belas agar harmonisasi aslinya dapat diukur sebagai acuan. Bila dua \textit{chorale} memakai melodi yang persis sama, yaitu interval dan ritme yang sama pada tinggi nada mana pun, hanya yang pertama dalam urutan korpus yang dapat terpilih, agar satu melodi tidak terhitung dua kali. Pemeriksaan program terhadap 350 berkas tersebut menemukan 23 berkas yang tidak memenuhi kriteria, terutama karena memiliki lebih dari empat suara, dan 9 berkas yang mengulang melodi \textit{chorale} lain, sehingga kerangka pengambilan sampel kelompok Bach berisi 318 melodi sebelum pemeriksaan kosakata DeepBach.

Seluruh latihan dalam buku Strube diinventaris terlebih dahulu dengan mencatat halaman, nomor, suara yang diberikan, tangga nada, birama, panjang, dan alasan kelayakan. Latihan bas tidak digunakan karena tugas penelitian ini memberikan sopran. Latihan melodi sesudah bab mode minor tidak digunakan karena dirancang untuk melatih nada nonakor, yaitu suspensi, nada sisipan, dan antisipasi. Pengecualian kaidah akibat nada nonakor memerlukan analisis akor. Pada melodi semacam itu, kejadian yang ditandai program tidak dapat dibedakan dari pelanggaran yang sebenarnya. Melodi pada bab harmonisasi \textit{chorale} juga tidak digunakan karena diambil dari J.S. Bach \parencite[174]{strube1928}. Inventaris latihan 1--118 (hlm. 8--81) menemukan 71 latihan yang memberikan melodi sopran pada bab yang ditetapkan, dan kelayakan setiap melodi diperiksa setelah melodi diketik. Seluruh latihan melodi yang layak diikutsertakan. Artinya, kelompok Strube adalah sensus atas latihan semacam itu dalam buku Strube, bukan sampel acak. Dalam pengujian statistik, melodi tersebut diperlakukan sebagai wakil soal latihan harmonisasi melodi dalam buku ajar harmoni, sehingga generalisasi di luar buku Strube dinyatakan secara terbatas. Melodi Bach yang layak kemudian diambil secara acak sebanyak jumlah melodi kelompok Strube, dengan jumlah sedikitnya 30 melodi. \textit{Seed} pengambilan sampel ditetapkan dalam protokol sebelum pengambilan dilakukan.

Jumlah minimum melodi ditetapkan melalui analisis daya statistik untuk uji Mann--Whitney pada pertanyaan ketiga. Untuk mendeteksi efek besar ($d = 0{,}8$) menurut Cohen \parencite{cohen1988} dengan taraf signifikansi 0,05 dua arah dan daya 0,8, uji \textit{t} dua kelompok memerlukan 25,5 melodi per kelompok. Efisiensi relatif asimtotik uji peringkat terhadap uji \textit{t} tidak kurang dari 0,864 untuk distribusi kontinu apa pun \parencite{hodges1956}, sehingga 25,5 dibagi 0,864 menjadi 29,5 dan dibulatkan ke atas menjadi 30 melodi per kelompok. Batas efisiensi tersebut berlaku untuk distribusi kontinu, sedangkan banyak nilai pelanggaran per birama mungkin nol, dan koreksi Holm pada analisis data menurunkan daya untuk sebagian pengujian. Karena itu, jumlah 30 melodi adalah target perencanaan, bukan jaminan daya statistik, dan ukuran efek terkecil yang dapat dideteksi dilaporkan berdasarkan jumlah melodi yang benar-benar dianalisis. Apabila latihan Strube yang layak kurang dari 30, seluruhnya tetap digunakan. Pertanyaan kedua menggunakan seluruh melodi sebagai pasangan, karena setiap melodi diharmonisasi oleh kedua model.

Setiap model menghasilkan lima harmonisasi untuk setiap melodi. Pengulangan ini mengurangi pengaruh hasil acak satu kali generasi terhadap nilai setiap melodi, dengan beban komputasi yang tetap terjangkau. Jumlah harmonisasi adalah jumlah melodi dikalikan dua model dan lima generasi, sehingga sedikitnya 600 harmonisasi untuk dua kelompok berisi 30 melodi. Kelima harmonisasi dari melodi yang sama tidak diperlakukan sebagai lima observasi independen; nilai pelanggaran per biramanya dirata-ratakan menjadi satu nilai per melodi per model. Setiap percobaan generasi, kegagalan, dan alasan eksklusi dicatat, dan keluaran yang gagal tidak diganti secara diam-diam.

Melodi \textit{chorale} Bach dapat termasuk data latih model. Menurut kode resmi DeepBach, contoh latihannya disusun mengikuti urutan korpus lalu dibagi berurutan tanpa diacak, yaitu 85 persen pertama untuk pelatihan, 10 persen berikutnya untuk validasi, dan 5 persen terakhir untuk evaluasi. Program penelitian ini menghitung ulang pembagian tersebut untuk setiap \textit{chorale}. Hasil hitungan dianggap sah hanya bila jumlah contoh latihannya sama persis dengan jumlah contoh dalam data yang disertakan bersama bobot DeepBach. Pembagian data latih \textit{checkpoint} Coconet yang digunakan tidak didokumentasikan, sehingga statusnya dicatat sebagai tidak diketahui. Dengan demikian, kelompok Bach dipahami sebagai melodi dari jenis musik yang dipelajari model, yang sebagian atau seluruhnya mungkin pernah dilihat model saat pelatihan.

\section{Metode Pengumpulan Data}\label{sec:metode-pengumpulan-data}

Data dikumpulkan melalui studi dokumen, generasi harmonisasi, dan pengukuran otomatis. Melodi dikumpulkan melalui studi dokumen, yaitu pengambilan melodi dari korpus \textit{chorale} dan dari buku Strube. Harmonisasi diperoleh dengan menjalankan kedua model pada melodi tersebut. Data penelitian, yaitu jumlah pelanggaran dan pelanggaran per birama, diperoleh melalui pengukuran otomatis terhadap harmonisasi dengan instrumen penelitian. Setiap tahap ditulis ke folder tersendiri yang tidak ditimpa, berisi identitas model, \textit{checkpoint}, versi perangkat lunak, pengaturan generasi, melodi, dan keluaran mentah setiap percobaan, agar pengukuran dapat diulang dan diperiksa.

\subsection{Variabel Penelitian}\label{sec:variabel-penelitian}

Penelitian ini memiliki dua variabel bebas dan satu variabel terikat.

\begin{enumerate}
    \item Variabel bebas pertama ($X_1$), yaitu model generatif. Variabel kategorik dengan dua kategori, DeepBach (jaringan rekuren LSTM dengan \textit{pseudo-Gibbs sampling}) dan Coconet (jaringan konvolusional dengan \textit{orderless NADE} dan \textit{blocked Gibbs sampling}). Kategori ini menunjukkan model dan \textit{checkpoint} yang diuji, bukan arsitektur yang terisolasi dari data latih dan prosedur generasi.
    \item Variabel bebas kedua ($X_2$), yaitu asal melodi. Variabel kategorik dengan dua kategori, melodi \textit{chorale} Bach dan melodi latihan Strube.
    \item Variabel terikat ($Y$), yaitu tingkat kepatuhan harmonisasi terhadap kaidah gerak suara Strube. Indikatornya adalah pelanggaran per birama untuk lima kaidah, yaitu kuint sejajar, oktaf sejajar, jarak antarsuara atas, persilangan di atas sopran, dan tumpang tindih suara. Kelima indikator dianalisis secara terpisah. Indeks penalti tertimbang menurut bobot rubrik Yan et al.\ \parencite{yan2018} digunakan sebagai ringkasan tambahan. Nilai yang lebih rendah menunjukkan kepatuhan yang lebih tinggi.
\end{enumerate}

Pertanyaan penelitian kedua membandingkan $Y$ antarkategori $X_1$ pada melodi yang sama. Pertanyaan penelitian ketiga menguji pengaruh $X_2$ terhadap $Y$ pada setiap kategori $X_1$.

\subsection{Pengumpulan Melodi}\label{sec:pengumpulan-melodi}

Melodi Bach diambil dari bagian sopran \textit{chorale} dalam korpus music21 beserta tangga nada, birama, dan tanda fermatanya, lalu disimpan dalam format MusicXML. Melodi Strube diketik peneliti dari edisi buku yang digunakan dalam format teks sederhana yang mencatat tanda kunci, mode, birama, birama gantung, serta nada, durasi, dan fermata setiap not. Program penelitian mengubah teks tersebut menjadi MusicXML dan membacanya kembali dengan pembaca yang sama yang dipakai instrumen, sehingga kesalahan pengetikan seperti jumlah ketukan yang tidak sesuai birama langsung ditolak. Setiap melodi diberi kode yang mencatat halaman dan nomor latihannya, dan hasil pengetikan dicocokkan kembali dengan cetakan. Hasil pengetikan dan inventaris disimpan sebagai data penelitian.

Untuk menggambarkan perbedaan kedua kelompok, ciri setiap melodi dihitung secara otomatis, yaitu jumlah birama, jumlah not, nada tertinggi dan terendah, lompatan terbesar, birama, mode, jumlah fermata, serta kesesuaian dengan batas Huang et al. Batas tersebut terpenuhi bila seluruh nada berada pada nada MIDI 60 sampai 81 dan tidak ada lompatan yang melebihi satu oktaf \parencite[798]{huang2019}.

\subsection{Generasi Harmonisasi}\label{sec:generasi-harmonisasi}

Setiap melodi diberikan kepada model sebagai sopran yang tetap, dan model menghasilkan alto, tenor, serta bas. Pada DeepBach, sopran ditetapkan sebagai batasan posisi, sedangkan informasi tangga nada, posisi ketukan, dan fermata diambil dari notasi melodi. Pada Coconet, sopran diisikan ke dalam \textit{piano roll}, dan model diminta mengisi seluruh langkah waktu alto, tenor, dan bas. Kedua model dijalankan dengan pengaturan bawaan implementasinya, yaitu 500 iterasi dan \textit{temperature} 1,0 untuk DeepBach serta 96 iterasi dan \textit{temperature} 0,99 untuk Coconet. \textit{Temperature} menentukan seberapa acak model memilih nada. Pengaturan tersebut dicatat dan sama untuk seluruh melodi. Setiap model dijalankan lima kali untuk setiap melodi. \textit{Seed} setiap percobaan DeepBach diturunkan dari satu \textit{seed} dalam protokol dan dicatat. Magenta.js tidak menyediakan pengaturan \textit{seed} untuk Coconet. Sebagai gantinya, seluruh keluaran mentah Coconet disimpan agar pengukurannya dapat diulang.

Setiap percobaan dicatat bersama statusnya, yaitu berhasil, ditolak karena melodi tidak dapat dibaca model, gagal karena galat model, atau gagal karena keluaran tidak dapat diubah menjadi empat suara. Percobaan yang gagal tidak dihapus dan tidak diulang secara diam-diam; pengulangan hanya dilakukan atas perintah tertulis dan dicatat sebagai percobaan baru. Keluaran kedua model diubah ke format MusicXML empat suara yang sama sebelum diukur. Keluaran Coconet berupa \textit{piano roll} yang hanya mencatat nada yang berbunyi pada setiap langkah seperenam belas, sehingga nada sama yang diulang tidak dapat dibedakan dari nada yang ditahan. Karena itu, nada sama yang berurutan pada alto, tenor, dan bas Coconet dibaca sebagai satu nada. Perbedaan representasi ini ditangani melalui analisis sensitivitas yang menyatukan nada sama yang berurutan pada setiap suara.

Sebelum generasi utama, uji coba (\textit{pilot}) dijalankan pada tiga melodi Bach yang tidak termasuk sampel utama dan tiga melodi Strube. Uji coba bertujuan menemukan masalah teknis, seperti panjang melodi yang tidak dapat diproses atau waktu generasi yang terlalu lama. Hasil uji coba disimpan terpisah dan tidak dipakai untuk menjawab pertanyaan penelitian.

\subsection{Instrumen Pengukuran}\label{sec:instrumen-pengukuran}

Instrumen evaluasi dikembangkan menggunakan bahasa Python dan pustaka music21 \parencite{cuthbert2010}. Kaidah yang diukur beserta pengecualiannya bersumber dari Strube \parencite{strube1928}, dan kategorinya mengikuti rubrik Yan et al.\ \parencite{yan2018}. Rumus pengukuran tidak dikutip dari kedua sumber tersebut. Peneliti menyusunnya sebagai definisi operasional, yaitu penerjemahan kaidah ke dalam kondisi yang dapat dihitung pada data simbolik. Halaman, contoh, dan pengecualian setiap kaidah pada edisi 1928 dicatat dalam lembar definisi kaidah, dan ringkasannya ditunjukkan pada Tabel~\ref{tab:kaidah}.

Keluaran diurai menjadi empat suara, yaitu sopran, alto, tenor, dan bas. Waktu dibagi ke dalam satuan not seperenam belas, yaitu resolusi waktu kedua model; Yan et al.\ juga menggunakan resolusi yang sama \parencite{yan2018}. Pada setiap satuan waktu, nada yang sedang berbunyi pada setiap suara dicatat, termasuk nada yang ditahan. Untuk setiap pasangan suara, peristiwa bunyi adalah satuan waktu tempat sedikitnya satu dari kedua suara tersebut memulai nada baru. Dalam rumus berikut, $t$ menunjukkan urutan peristiwa bunyi pasangan yang diperiksa, $p_{v,t}$ adalah nomor nada MIDI suara $v$ pada peristiwa $t$, dengan $v = 1, 2, 3, 4$ untuk sopran, alto, tenor, dan bas, dan $\Delta p_{v,t} = p_{v,t} - p_{v,t-1}$. Untuk dua suara $i < j$, $d_{ij,t} = p_{i,t} - p_{j,t}$ adalah jarak suara atas terhadap suara bawah dalam semiton. Gerak antara dua peristiwa hanya diperiksa apabila kedua suara berbunyi pada kedua peristiwa tersebut.

Kuint sejajar dan oktaf sejajar diperiksa pada keenam pasangan suara. Untuk interval target $\delta = 7$ (kuint murni, termasuk kuint majemuk) dan $\delta = 0$ (oktaf, oktaf majemuk, dan unisono), kejadian dihitung bila kedua suara bergerak searah dari interval target ke interval target. Gerak berlawanan tidak dihitung karena Strube tidak menganggapnya bermasalah \parencite[9]{strube1928}, dan kuint murni ke kuint berkurang tidak dihitung karena kuint kedua tidak murni \parencite[35]{strube1928}.

{\small
\begin{multline}
\text{Paralel}_{ij}^{\delta}(t) = \mathds{1}\left(|d_{ij,t-1}| \bmod 12 = \delta\right) \times \mathds{1}\left(|d_{ij,t}| \bmod 12 = \delta\right) \\
\times \mathds{1}\left(\Delta p_{i,t} \neq 0\right) \times \mathds{1}\left(\text{sgn}(\Delta p_{i,t}) = \text{sgn}(\Delta p_{j,t})\right)
\end{multline}
}

Jarak antarsuara atas diperiksa pada pasangan sopran--alto dan alto--tenor \parencite[20]{strube1928}. Persilangan diperiksa antara sopran dan setiap suara lainnya, yaitu sopran--alto, sopran--tenor, dan sopran--bas, karena Strube membolehkan persilangan sesekali kecuali yang melewati sopran \parencite[174]{strube1928}.

{\small
\begin{equation}
\text{Jarak}_{ij}(t) = \mathds{1}\left(d_{ij,t} > 12\right), \qquad \text{Silang}_{ij}(t) = \mathds{1}\left(d_{ij,t} < 0\right)
\end{equation}
}

Tumpang tindih suara diperiksa pada tiga pasangan suara yang bersebelahan. Kejadian dihitung apabila suara bawah bergerak lebih tinggi daripada nada sebelumnya pada suara atas, atau suara atas bergerak lebih rendah daripada nada sebelumnya pada suara bawah \parencite[12]{strube1928}. Kejadian tersebut hanya dihitung bila kedua suara tidak bersilangan pada peristiwa tersebut, agar satu kejadian tidak dihitung dua kali, dan bila tidak ada suara pada pasangan itu yang bergerak melangkah, yaitu satu atau dua semiton, karena Strube membolehkan tumpang tindih dengan gerak melangkah \parencite[12--13]{strube1928}. Dengan $\text{Langkah}_{v}(t) = \mathds{1}\left(|\Delta p_{v,t}| \in \{1, 2\}\right)$, kejadian tumpang tindih adalah

{\small
\begin{multline}
\text{Tumpang}_{ij}(t) = \mathds{1}\left(p_{j,t} > p_{i,t-1} \lor p_{i,t} < p_{j,t-1}\right) \times \mathds{1}\left(d_{ij,t} \geq 0\right) \\
\times \left(1 - \max\left(\text{Langkah}_{i}(t), \text{Langkah}_{j}(t)\right)\right)
\end{multline}
}

Karena sopran diberikan sebagai soal, setiap pasangan yang diperiksa melibatkan sedikitnya satu suara yang dihasilkan model. Dengan demikian, setiap kejadian yang ditandai mencerminkan keputusan model atas melodi yang diberikan. Ringkasan setiap kaidah ditunjukkan pada Tabel~\ref{tab:kaidah}.

\begin{table}[ht]
\centering
\small
\begin{tabularx}{\textwidth}{|l|l|>{\raggedright\arraybackslash}X|c|c|}
\hline
\textbf{Kaidah} & \textbf{Pasangan} & \textbf{Kejadian yang dihitung} & \textbf{Strube} & \textbf{Bobot Yan} \\ \hline
Kuint sejajar & Keenam pasangan & Gerak searah dari kuint murni ke kuint murni & 9, 12 & 1 \\ \hline
Oktaf sejajar & Keenam pasangan & Gerak searah dari oktaf atau unisono ke oktaf atau unisono & 9, 12 & 1 \\ \hline
Jarak suara atas & S--A, A--T & Peristiwa bunyi dengan jarak lebih dari satu oktaf & 20 & 0,5 \\ \hline
Persilangan & S--A, S--T, S--B & Peristiwa bunyi dengan suara lain di atas sopran & 174 & 0,5 \\ \hline
Tumpang tindih & S--A, A--T, T--B & Gerak melewati nada sebelumnya pada suara sebelah tanpa gerak melangkah & 12--13 & 0,5 \\ \hline
\end{tabularx}
\caption{Kaidah gerak suara yang diukur, halaman uraiannya dalam Strube (1928), dan bobotnya dalam rubrik Yan et al.\ (2018)}
\label{tab:kaidah}
\end{table}

Pelanggaran per birama untuk kaidah $r$ pada satu harmonisasi adalah $L_r = N_r / M$, dengan $N_r$ jumlah kejadian kaidah tersebut dan $M$ jumlah birama melodi masukan; birama gantung dihitung sebagai satu birama. Satuan per birama mengikuti Huang et al.\ \parencite[798]{huang2019}. Indeks penalti tertimbang adalah

{\small
\begin{equation}
P = \frac{N_{\text{kuint}} + N_{\text{oktaf}} + 0{,}5 \left(N_{\text{jarak}} + N_{\text{silang}} + N_{\text{tumpang}}\right)}{M}
\end{equation}
}

Bobot pada indeks tersebut mengikuti rubrik Yan et al.\ \parencite{yan2018}, tetapi indeksnya sendiri adalah adaptasi peneliti, bukan skor rubrik Yan et al. Indeks ini meringkas kelima kaidah, sedangkan pengujian hipotesis dilakukan pada setiap kaidah. Jumlah kejadian dan jumlah kesempatan, yaitu banyaknya gerak atau peristiwa bunyi yang diperiksa, juga disimpan untuk setiap kaidah.

Penghitungan utama mencakup semua kejadian, mengikuti Huang et al. Dua penghitungan tambahan dilakukan untuk analisis sensitivitas. Penghitungan pertama mengecualikan gerak yang berawal ketika sopran sedang membunyikan nada berfermata, karena gerak paralel pada batas frasa dapat dimaklumi \parencite[798]{huang2019} dan Strube memperlakukan fermata sebagai kadens \parencite[174]{strube1928}. Pengecualian ini mencakup seluruh durasi nada berfermata, bukan hanya perpindahan ke frasa berikutnya. Penghitungan kedua menyatukan nada sama yang berurutan pada setiap suara sebelum pengukuran, sehingga keluaran DeepBach, keluaran Coconet, dan harmonisasi Bach dibaca dengan representasi bunyi yang sama. Pengecualian akibat pengulangan akor dan nada nonakor \parencite[35]{strube1928} tidak diterapkan karena memerlukan analisis akor.

\subsection{Validitas dan Reliabilitas Instrumen}\label{sec:validitas-dan-reliabilitas}

Instrumen yang valid mengukur secara akurat apa yang hendak diukur, sedangkan instrumen yang reliabel memberikan hasil yang konsisten ketika digunakan berulang kali pada objek yang sama \parencite{leedy2018}. Lulus uji perangkat lunak hanya berarti program berperilaku benar pada kasus uji, sedangkan validitas menyangkut apakah kejadian yang ditandai sesuai dengan kaidah yang dimaksud. Validitas instrumen diperiksa melalui empat langkah yang tidak memerlukan penilai manusia.

\begin{enumerate}
    \item Validitas isi. Setiap kaidah dikaitkan dengan halaman uraian Strube dan kategori rubrik Yan et al. Definisi, contoh, dan pengecualian ditetapkan sebelum harmonisasi utama dihasilkan.
    \item Uji terhadap contoh buku. Contoh gerak suara yang benar dan salah yang dicetak dalam buku Strube diketik sebagai kasus uji. Instrumen harus menandai kesalahan yang ditunjukkan buku dan tidak menandai contoh yang benar. Kasus batas yang disusun peneliti, seperti tanda diam, nada yang ditahan, interval majemuk, unisono, serta perbedaan persilangan dan tumpang tindih, ditambahkan. Setiap ketidaksesuaian dilaporkan.
    \item Pemeriksaan silang. Pemeriksaan dasar gerak paralel, persilangan, dan tumpang tindih pada seluruh \textit{chorale} Bach dalam korpus music21 dibandingkan dengan fungsi \texttt{VoiceLeadingQuartet} dari music21 yang memeriksa gerak yang sama. Fungsi paralel music21 juga menerima kuint dan oktaf yang dicapai dengan gerak berlawanan, yang tidak dianggap bermasalah oleh Strube, sehingga perbedaan jenis itu dipisahkan. Kesesuaian ini menunjukkan ketepatan program pada pemeriksaan dasar, bukan validitas semua kejadian yang ditandai.
    \item Perbandingan dengan angka terbitan. Kuint dan oktaf sejajar per birama pada \textit{chorale} Bach dibandingkan dengan angka Huang et al., yaitu 0,023 dan 0,009 \parencite[798]{huang2019}. Angkanya tidak diharapkan sama persis karena versi data dan rincian definisi dapat berbeda. Perbandingan ini hanya pemeriksaan kewajaran.
\end{enumerate}

Reliabilitas instrumen diperiksa melalui pengulangan pengukuran pada berkas yang sama dengan versi instrumen yang sama. Karena instrumen tidak mengandung unsur acak, pengukuran yang reliabel menandai kejadian yang sama persis. Variasi antarharmonisasi berasal dari proses generasi model, dan variasi tersebut ditangani melalui lima harmonisasi per melodi. Perubahan satuan waktu, pemetaan suara, definisi kaidah, atau penyebut (jumlah birama) diperlakukan sebagai versi instrumen baru, dan hasil versi sebelumnya disimpan.

Kejadian yang ditandai instrumen adalah pelanggaran menurut definisi operasional. Instrumen tidak mempertimbangkan pengecualian yang memerlukan penafsiran akor, sehingga hasil penelitian tidak dinyatakan sebagai penilaian musikal atas setiap kejadian. Strube sendiri membolehkan sopran sesekali berjarak lebih dari satu oktaf dari alto \parencite[20]{strube1928}, sehingga tanda pada kaidah jarak berarti jaraknya melewati ambang, belum tentu kesalahan.

\subsection{Kontrol Kualitas Sampel}\label{sec:kontrol-kualitas-sampel}

Harmonisasi layak dianalisis bila memuat empat suara yang dapat dipisahkan, setiap suara hanya membunyikan satu nada pada satu waktu, seluruh nada dapat dinyatakan dalam satuan not seperenam belas, panjang dan sopran harmonisasi identik dengan melodi masukan dalam tinggi nada dan saat mulainya, dan setiap suara yang dihasilkan model memuat sedikitnya satu nada. Tanda diam di dalam suara tetap diperbolehkan. Pemetaan suara perlu diperiksa karena pemisahan suara pada musik simbolik merupakan persoalan tersendiri, terutama ketika suara bersilangan atau tumpang tindih \parencite{karystinaios2023}. Harmonisasi yang tidak layak dikeluarkan dari analisis utama, dan jumlah serta alasan kegagalannya dilaporkan per model dan per asal melodi. Kegagalan itu juga termasuk hasil penelitian dan tidak dibuang begitu saja. Melodi yang tidak memiliki harmonisasi layak dari salah satu model dikeluarkan dari perbandingan berpasangan untuk model tersebut, dan pengeluaran itu dicatat.

\section{Metode Analisis Data}\label{sec:metode-analisis-data}

Nilai pelanggaran per birama setiap kaidah dirata-ratakan dari harmonisasi yang layak untuk setiap melodi dan setiap model. Nilai rata-rata tersebut menjadi data analisis, dan banyaknya harmonisasi layak dari lima percobaan dicatat untuk setiap melodi dan model.

Pertanyaan penelitian pertama dijawab dengan statistik deskriptif. Statistik tersebut mencakup jumlah percobaan generasi, harmonisasi yang layak, dan kegagalan pada setiap model dan asal melodi; median, rentang antarkuartil, rata-rata, dan deviasi standar pelanggaran per birama setiap kaidah; serta nilai harmonisasi asli Bach pada melodi kelompok Bach sebagai acuan. Selisih nilai setiap model terhadap harmonisasi Bach pada melodi yang sama juga dilaporkan. Ciri melodi kedua kelompok, termasuk proporsi melodi yang berada di luar batas Huang et al., dilaporkan untuk menunjukkan seberapa jauh kedua kelompok berbeda.

Pertanyaan penelitian kedua diuji dengan uji peringkat bertanda Wilcoxon dua arah pada selisih nilai DeepBach dikurangi Coconet untuk setiap melodi, karena setiap melodi diharmonisasi oleh kedua model sehingga datanya berpasangan. Selisih bernilai nol tidak diikutsertakan dalam perhitungan peringkat, dan bila seluruh selisih bernilai nol, hasilnya dilaporkan tanpa nilai $p$. Pertanyaan penelitian ketiga diuji dengan uji Mann--Whitney untuk setiap model, karena kelompok Bach dan kelompok Strube terdiri atas melodi yang berbeda. Uji Mann--Whitney juga digunakan Huang et al.\ untuk membandingkan kelompok melodi pada analisis Bach Doodle \parencite[798]{huang2019}. Uji untuk kuint dan oktaf sejajar pada pertanyaan ketiga dilakukan satu arah sesuai hipotesis, dengan nilai $p$ dua arah dilaporkan juga, sedangkan uji lainnya dilakukan dua arah. Uji nonparametrik dipilih karena pelanggaran per birama tidak dapat bernilai negatif, banyak bernilai nol, dan tidak dapat diasumsikan berdistribusi normal. Ukuran efek dilaporkan sebagai korelasi peringkat biserial; nilai positif berarti pelanggaran per birama lebih tinggi pada DeepBach (pertanyaan kedua) dan pada melodi Strube (pertanyaan ketiga). Taraf signifikansi ditetapkan 0,05, dan koreksi Holm diterapkan pada lima kaidah dalam setiap kelompok pengujian, yaitu satu kelompok untuk pertanyaan kedua dan satu kelompok untuk setiap model pada pertanyaan ketiga. Indeks penalti tertimbang dilaporkan secara deskriptif dan diuji dengan cara yang sama sebagai analisis tambahan di luar koreksi tersebut.

Analisis sensitivitas berikut dilakukan untuk melihat apakah kesimpulan berubah bila cara penghitungan diubah:

\begin{enumerate}
    \item pengujian diulang dengan penghitungan yang mengecualikan gerak di sekitar fermata;
    \item pada DeepBach, melodi Bach yang termasuk data latih dibandingkan dengan yang tidak termasuk, hanya bila pembagian data latihnya telah dinyatakan sah;
    \item pengujian diulang hanya pada melodi yang kelima harmonisasinya layak;
    \item pengujian diulang dengan penghitungan yang menyatukan nada sama yang berurutan.
\end{enumerate}

Hasil yang tidak signifikan tidak ditafsirkan sebagai bukti bahwa model atau asal melodi setara. Seluruh keluaran diukur dengan instrumen dan prosedur yang sama. Ukuran efek terkecil yang dapat dideteksi dilaporkan berdasarkan jumlah melodi yang benar-benar dianalisis.

Hasil disajikan dalam bentuk tabel dan diagram kotak (\textit{box plot}) pelanggaran per birama setiap kaidah per model dan per asal melodi. Contoh kejadian yang ditandai untuk pembahasan dipilih secara acak dengan \textit{seed} yang dicatat, lalu ditampilkan beserta nomor birama dan kutipan partiturnya. Pembahasan contoh tersebut menjelaskan konteks musikal setiap kejadian, termasuk kemungkinan adanya nada nonakor atau batas frasa, tanpa mengubah data hasil pengukuran.

\section{Rancangan Alur Penelitian}\label{sec:rancangan-alur-penelitian}

Penelitian berlangsung dalam tahap berikut.

\begin{enumerate}
    \item Perumusan kaidah dan inventaris melodi, yaitu penetapan edisi dan halaman kaidah Strube, contoh, pengecualian, dan definisi operasional; inventaris latihan Strube; serta pengetikan melodi yang layak.
    \item Pengembangan instrumen, yaitu penulisan dan pengujian program evaluasi berbasis music21.
    \item Validasi instrumen, yaitu uji terhadap contoh buku Strube, pemeriksaan silang dengan fungsi music21, dan perbandingan dengan angka Huang et al.
    \item Penyiapan melodi, yaitu pengubahan melodi Bach dan Strube ke MusicXML, pemeriksaan kelayakan, dan pemeriksaan masukan terhadap kedua model tanpa menjalankan model.
    \item Uji coba (\textit{pilot}), yaitu generasi dan pemeriksaan sejumlah kecil harmonisasi dari setiap model dan asal melodi untuk menemukan masalah teknis sebelum protokol ditetapkan. Hasil uji coba dipisahkan dari data utama.
    \item Penetapan protokol dan pemilihan sampel, yaitu pembekuan definisi, pengaturan generasi, dan rencana analisis, lalu pengambilan acak melodi Bach.
    \item Pengumpulan data, yaitu generasi lima harmonisasi per melodi untuk setiap model.
    \item Evaluasi otomatis, yaitu pemeriksaan kualitas dan pengukuran seluruh harmonisasi, dengan hasil disimpan dalam format CSV terstruktur.
    \item Analisis data, yaitu statistik deskriptif, uji hipotesis, analisis sensitivitas, dan pembahasan contoh.
\end{enumerate}

\begin{figure}[htbp]
\centering
\begin{tikzpicture}[node distance=1.35cm, every node/.style={fill=white, font=\footnotesize}, align=center]
  \tikzset{
    startstop/.style={rectangle, rounded corners, minimum width=3.5cm, minimum height=0.7cm, text centered, draw=black},
    process/.style={rectangle, minimum width=3.5cm, minimum height=0.7cm, text centered, draw=black},
    decision/.style={diamond, aspect=2, minimum width=3.5cm, minimum height=0.7cm, text centered, draw=black},
    arrow/.style={thick,->,>=stealth}
  }

  % Nodes
  \node (start) [startstop] {Mulai};
  \node (rule) [process, below of=start] {Perumusan Kaidah Strube \\ \& Inventaris Melodi};
  \node (dev) [process, below of=rule] {Pengembangan Program \\ Evaluasi (music21)};
  \node (val) [process, below of=dev] {Validasi Instrumen \\ (Contoh Strube, music21, Huang)};
  \node (dec) [decision, below of=val, yshift=-0.1cm] {Apakah Valid?};
  \node (prep) [process, below of=dec, yshift=-0.3cm] {Penyiapan Melodi \\ \& Pemeriksaan Masukan};
  \node (pilot) [process, below of=prep] {Uji Coba \\ \& Penetapan Protokol};
  \node (gen) [process, below of=pilot] {Generasi Harmonisasi};
  \node (eval) [process, below of=gen] {Evaluasi Otomatis \\ \& Simpan CSV};
  \node (anal) [process, below of=eval] {Analisis Data \\ \& Uji Hipotesis};
  \node (stop) [startstop, below of=anal] {Selesai};

  % Connections
  \draw [arrow] (start) -- (rule);
  \draw [arrow] (rule) -- (dev);
  \draw [arrow] (dev) -- (val);
  \draw [arrow] (val) -- (dec);
  \draw [arrow] (dec) -- node[anchor=east] {Ya} (prep);
  \draw [arrow] (prep) -- (pilot);
  \draw [arrow] (pilot) -- (gen);
  \draw [arrow] (gen) -- (eval);
  \draw [arrow] (eval) -- (anal);
  \draw [arrow] (anal) -- (stop);

  % Feedback loop
  \draw [arrow] (dec.east) -- ++(1.2,0) |- node[anchor=south, pos=0.25] {Tidak} (dev.east);

\end{tikzpicture}
\caption{Rancangan Alur Penelitian}
\label{fig:diagram-alir}
\end{figure}
